% Nouveau EGA V (cote 25): archive cote 25, site page 129 (PDF page 130).
% English translation of the checked French annotated master.
% AI-translated by GPT-6 Astra (OpenAI, ChatGPT Pro mode).
\scanpage{25/129}
\mspage{\hand{32}}
% NOTE: A pencil oblique stroke follows the upper page number.

preceding, and from the determination of the automorphisms of $\overset{\hand{\vee}}{\mathrm{P}}(E) = P(\overset{\hand{\vee}}{E})$ made
in no. 1 , which implies that an automorphism of a projective bundle over an S
(relative to a Module loc free of finite type) is known when one knows
its effect on a family of sections forming a ``projective basis''
(i.e. forming a projective basis in each \var{\xout{point}}{}{typewriter strike-out}{ega25_25_129-v0001} fiber).

Let us now sketch the general case: arbitrary S. Of course,
by making a base change on S, one is reduced in a) to prov
that two sections of $\overset{\hand{\vee}}{\mathrm{P}}$ over S wh\uline{i}ch define the same divisor on X are
identical, and in b) that two automorphisms of $\overset{\hand{\vee}}{\mathrm{P}}$ such that
$\var{\over{D}{D}}{D}{writing over}{ega25_25_129-v0003}\circ u = D\circ v$ are identical. \var{\inserted[typed]{One can assume S affine,}}{One can assume S affine,}{insertion}{ega25_25_129-v0004} In case b), where g was not expressly assumed
to be of finite presentation, but was assumed to be surjective, one immediately reduces
(thanks to the fact that g is open) to the case where X is also affine, hence
of finite presentation over S. By the patented procedure, one can now
reduce to the case of \uline{noetherian} S. Now for a base prescheme \var{\inserted[typed]{S}}{S }{insertion}{ega25_25_129-v0006}\var{\inserted{loc}}{loc}{insertion}{ega25_25_129-v0007} noetherian,
and a morphism of functors $F \overset{\hand{h}}{\hand{\longrightarrow}} G$ over S (F and G functors
$(\mathrm{Sch})_{/S}^{\circ} \hand{\longrightarrow} (\mathrm{Ens})$ ), there is \var{\xout{some}}{}{typewriter strike-out}{ega25_25_129-v0008}\unsure{a} \var{\xout{criteria}}{criterion}{typewriter strike-out}{ega25_25_129-v0009} that is very \unsure{general} \var{\inserted[typed]{(which will be made explicit in Chap. V),}}{(which will be made explicit in Chap. V),}{insertion}{ega25_25_129-v0011} allowing one to assert
that when for every $s \hand{\in} S$, the corresponding morphism $F_s \hand{\longrightarrow} G_s$ is
a monomorphism, then $F \hand{\longrightarrow} G$ is a monomorphism (NB one puts $F_s{=}F\times_S\mathrm{Spec}$
$(k(s))$, and similarly for $G_s$ ), subject to simple hypotheses on F and G
(satisfied for example when F and G are both representable by
preschemes loc of finite type over S; but in the case before us, only
the first of the two functors is representable a priori). We shall
sketch the argument of Chap. V in the \var{\inserted[typed]{two}}{two}{insertion}{ega25_25_129-v0014} particular cases that
interest us here. There are two sections \var{\inserted{u, v}}{u, v}{insertion}{ega25_25_129-v0015} (of $\overset{\hand{\vee}}{\mathrm{P}}$, resp. of the projective group
\uline{GP} ($\overset{\hand{\vee}}{E}$) ) of a \hand{pre}scheme of finite type \var{\inserted{F}}{F}{insertion}{ega25_25_129-v0016} over S that one wants to show
are equal; for this it clearly suffices to prove that they are
equal after every base change $\mathrm{Spec}(\underline{O}_{S,s}/\underline{m}_s^{\hand{n+1}}$ ) $\hand{\longrightarrow}$ \hand{$\mathcal{S}$},

\mspage{\hand{129}}
\tnote{25-129-1}{The source ``à prouve'' is clipped at the edge; the English fragment ``to prov'' retains its incompleteness. The marked word ``general'' represents uncertainty about the final l of the French word. The inserted abbreviation loc means locally and qualifies noetherian; it remains at the source insertion point after the inserted base letter. The underlined i in ``which'' corresponds to the underlined i of French ``qui''. The final x canceled in ``deux'' leaves ``deu'' on the source page; this French letter cancellation has no separate English effect.}
